\begin{table}[H]
\centering
\caption{Standardized Effect Sizes for Main Outcomes}
\label{tab:sde}
\begin{adjustbox}{max width=\textwidth}
\begin{threeparttable}
\small
\begin{tabular}{lcccccc}
\toprule
Outcome & $\hat{\beta}$ & SE & SD($Y$) & SDE & SE(SDE) & Class. \\
\midrule
\multicolumn{7}{l}{\textit{Panel A: Pooled}} \\
Employment (all races) & 0.5152 & 0.2171 & 1.926 & 0.2675 & 0.1127 & Large positive \\
Employment (Asian interaction) & -0.0189 & 0.2980 & 1.926 & -0.0098 & 0.1547 & Small negative \\
Earnings (all races) & -0.1809 & 0.0745 & 0.436 & -0.4150 & 0.1709 & Large negative \\
\midrule
\multicolumn{7}{l}{\textit{Panel B: Heterogeneous (Asian interaction)}} \\
High-Asian-share industries & -0.1770 & 0.5771 & 1.964 & -0.0901 & 0.2939 & Moderate negative \\
Low-Asian-share industries & -0.0570 & 0.3844 & 1.916 & -0.0298 & 0.2007 & Small negative \\
\bottomrule
\end{tabular}
\begin{tablenotes}[flushleft]
\scriptsize
\item \textit{Notes:} \textbf{Country:} United States. \textbf{Research question:} Whether the 2018--2019 Section 301 tariffs on Chinese imports produced racially heterogeneous manufacturing employment effects, given the 4:1 overrepresentation of Asian workers in tariff-targeted electronics manufacturing. \textbf{Policy mechanism:} Section 301 tariffs imposed ad valorem duties of 10--25\% on Chinese imports across manufacturing industries, raising input costs and reducing demand for import-competing products, with the heaviest rates on electronics (NAICS 334) and machinery (NAICS 333) where Asian workers are disproportionately employed. \textbf{Outcome definition:} Log beginning-of-quarter county-industry-race employment (Emp) from the QWI, and log average monthly stable earnings (EarnS). \textbf{Treatment:} Continuous---Bartik county-level tariff exposure constructed as the employment-weighted sum of industry-level Section 301 tariff rates using 2017Q2 baseline shares. \textbf{Data:} QWI race $\times$ 3-digit NAICS $\times$ county $\times$ quarter panels (Census Bureau), 2015Q1--2019Q4; Chinese import penetration from Census International Trade; Section 301 rates from USTR. \textbf{Method:} OLS with county$\times$race, industry$\times$quarter, and race$\times$quarter fixed effects; SEs clustered at the state level. \textbf{Sample:} U.S.\ manufacturing industries (NAICS 31--33), county-industry-race-quarter cells with positive employment; three race groups (White, Black, Asian). SDE $= \hat{\beta} / \text{SD}(Y)$ where SD($Y$) is the pre-treatment standard deviation. Classification refers to magnitude, not statistical significance: Large ($|$SDE$| > 0.15$), Moderate ($0.05$--$0.15$), Small ($0.005$--$0.05$), Null ($< 0.005$).
\end{tablenotes}
\end{threeparttable}
\end{adjustbox}
\end{table}

